\section{问题分析}

\subsection{总体思路}

若把四个问题分别建立四套模型，容易在符号、储能效率、时间索引和费用口径上产生不一致。实际上，四问中的微网物理结构始终不变：负载需要被满足，光伏优先在本地消纳，储能受容量、功率和效率约束，外网负责补足剩余需求。真正变化的只有两个问题：\textbf{决策时刻知道什么}，以及\textbf{计划发生偏差后怎样付费}。因此本文把能量平衡和储能递推作为公共内核，再逐问改变信息集、目标函数和结算规则。

这一处理使四问形成自然的递进关系。问题一先在确定信息下回答“若未来已知，储能应怎样与外网购电协同”；问题二取消当天真实信息，回答“预测误差存在时，计划要保守到什么程度，以及储能能缓冲多少偏差”；问题三引入日内新预报，回答“较新的信息值多少钱、应不应该重新调整计划”；问题四进一步改变价格环境，检验上述策略在实时价格波动下是否仍然成立。整篇论文的求解顺序也据此安排为“数据与时间口径统一—公共物理模型—分问信息与费用模型—结果解释—稳健性检验”。

\subsection{各问题的建模重点}

\indent\textbf{问题一：}附件1已给定典型日的电价、负载和光伏预测，可以直接进行确定性经济调度。储能的作用来自两个方向：一方面在低价时段买电充电、在高价时段放电；另一方面吸收中午可能出现的光伏富余。题目又要求 $0{:}00$ 与 $24{:}00$ 储电量相同，因此不能通过“把初始电量在一天内耗尽”人为降低费用，首末状态必须作为硬约束。

\indent\textbf{问题二：}与问题一相比，最大的变化不是物理设备，而是全天计划必须在真实源荷完全实现之前确定。本文采用严格因果口径：主模型只利用 $0{:}00$ 前已经发生的历史数据形成当天负载、光伏预测；若实际业务能够提前提供全天计划负荷，则把该情形作为理想信息基准，而不与主模型混用。由于紧急购电价格是正常价格的 $5$ 倍，高估光伏比低估光伏的经济后果更严重，故光伏预测不应只追求均方意义上的“无偏”，还需考虑非对称风险。计划执行后，外网全天计划不再改变，但储能可根据当前已实现偏差实时修正充放电，最后无法覆盖的缺口才进入紧急购电。

\indent\textbf{问题三：}本问的核心不再是“能否预测”，而是“新预报到来后是否值得改计划”。因此每日 $0{:}00$ 先给出覆盖全天的基准计划；在 $6{:}00$、$12{:}00$、$18{:}00$ 获得更新预报后，只重算尚未执行的时段，并把当前实测储电量作为新的初始状态。为了回答“其他时刻预报是否需要引入”，不能只展示某一天的曲线，而应保持其他条件不变，逐步增加可用更新节点，比较全年费用变化。由于节点更新时还会同时获得已经发生的负载信息，主对照反映的是节点的综合价值，后续再用配对控制区分负载刷新与新光伏预报的增量贡献。

\indent\textbf{问题四：}题面直接给出了附件4的全年实时电价，因此正式重算以附件4电价为已给数据，将其替换问题二、三中的典型日电价即可。这样可以单独观察价格波动对储能充放电时点和购电成本的影响。考虑真实市场中未来价格未必一次性全部公布，本文另设因果价格预测作为扩展，但不改变题面主结果。

\subsection{求解方法选择}

上述关系均可写成线性等式、不等式和分段线性费用。问题一、问题二以及不含充放电互斥整数变量的对照模型可用线性规划求解；需要严格禁止同一时段同时充、放电时，引入二元变量后形成混合整数线性规划。两类模型均使用 HiGHS 求解。相比元启发式算法，这一做法能够直接利用问题的凸结构，所得结果易于复算，也便于逐时检查能量守恒和费用组成。
